\documentclass[12pt]{article}
\usepackage[utf8]{inputenc}
\usepackage[T1]{fontenc}
\usepackage[margin=1in]{geometry}
\usepackage{setspace}
\usepackage{xcolor}
\usepackage[colorlinks=true, urlcolor=blue!60!black, linkcolor=blue!60!black]{hyperref}
\setlength{\parindent}{2em}
\setlength{\parskip}{1.25em}
\renewcommand{\baselinestretch}{1.0}
\begin{document}
\begin{center}
{\LARGE \textit{TransUnion LLC v. Ramirez}, 594 U.S. 413 (2021)\par}
\vspace{0.8em}
{\large Case Brief --- Four-Page Edition\par}
\end{center}
\vspace{1em}

\section{Memory Jogger}

After TransUnion falsely flagged 8,185 consumers as potential matches to a federal terrorist list in violation of the Fair Credit Reporting Act, the Court held 5--4 that only the 1,853 whose reports were actually sent to businesses suffered a ``concrete'' Article III injury; the rest lacked standing because ``an injury in law is not an injury in fact'' (\href{https://supreme.justia.com/cases/federal/us/594/20-297/}{Justia}).

\section{Detailed Case Facts}

The FCRA requires credit reporting agencies to ``follow reasonable procedures to assure maximum possible accuracy,'' to disclose on request ``[a]ll information in the consumer's file at the time of the request,'' and to provide a summary of rights with each written disclosure; willful violations carry statutory damages of \$100 to \$1,000 plus punitive damages. TransUnion's ``OFAC Name Screen Alert'' compared only consumers' first and last names with Treasury's list of terrorists, drug traffickers, and other serious criminals, producing thousands of false positives.

In February 2011 a Nissan dealer refused to sell Sergio Ramirez a car because his report showed the alert. When he requested his file, TransUnion sent one mailing that omitted the alert and a separate ``courtesy'' letter the next day noting the potential match but lacking the summary of rights. Alarmed, Ramirez canceled a planned trip to Mexico.

\section{Procedural History}

Ramirez sued for willful violations on behalf of 8,185 consumers. The parties stipulated that reports with OFAC alerts went to third parties for only 1,853 class members. The district court found standing for all; a jury awarded each member \$984.22 in statutory and \$6,353.08 in punitive damages, over \$60 million total. A divided Ninth Circuit affirmed standing but reduced punitive damages; Judge McKeown would have limited standing. The Court granted certiorari and decided the case on June 25, 2021.

\section{Judicial Votes}

Reversed and remanded, \textbf{5--4}.
\begin{itemize}
\item \textbf{Majority:} Justice Kavanaugh, joined by Chief Justice Roberts and Justices Alito, Gorsuch, and Barrett.
\item \textbf{Dissent:} Justice Thomas, joined by Justices Breyer, Sotomayor, and Kagan.
\item \textbf{Dissent:} Justice Kagan, joined by Justices Breyer and Sotomayor.
\end{itemize}

\section{Holding}

Only plaintiffs concretely harmed by a statutory violation may sue a private defendant for damages, and every class member must have standing to recover. An intangible harm is concrete if it bears a ``close relationship''---not an ``exact duplicate''---to a harm traditionally actionable in American courts. Dissemination of the misleading alert resembles defamation and confers standing; undisseminated internal files do not; an unmaterialized risk of future harm cannot support retrospective damages; and informational injuries without ``downstream consequences'' fail. Only Ramirez had standing on the mailing claims. The Court reserved whether standing must be shown before class certification.

\section{Analysis of Opinions}

\textbf{Majority (Kavanaugh).} Article III confines courts to real disputes, not a ``roving commission.'' Under \textit{Spokeo}, Congress may ``elevate'' real harms but ``may not simply enact an injury into existence''; letting uninjured plaintiffs sue would also intrude on the Executive's enforcement discretion. Publication is ``essential to liability'' in defamation, so an inaccurate file in a database is like a defamatory letter in a desk drawer. A risk of harm may support an injunction under \textit{Clapper}, but not damages for a risk that never materialized---like a reckless driver who misses a commuter who gets home safely. There was no evidence other members even opened the mailings, and \textit{Akins} and \textit{Public Citizen} involved public-disclosure laws.

\textbf{Thomas (dissenting).} Founding-era law distinguished private rights, whose violation alone supported suit, from public rights, which required actual damage. ``Injury in fact'' first appeared in 1970. Each class member's private statutory rights were violated, and even under the majority's test the class had standing; he warned the ruling may be a ``pyrrhic victory'' that sends such class actions to state courts.

\textbf{Kagan (dissenting).} The Court turned standing into a tool of ``judicial aggrandizement''; courts should override Congress's authorization to sue only when it could not reasonably have thought a suit would address real harm.

\section{Examples: Future Applications}

\begin{enumerate}
\item \textbf{Same-Side (No Standing): Retained Data.} A former cable subscriber whose data was kept past the statutory deadline, without breach or disclosure, alleges no concrete harm.
\item \textbf{Same-Side (No Standing): Small-Font Notice.} A debt-collection letter with dispute rights in the wrong font, causing no confusion, is a bare procedural violation.
\item \textbf{Opposite-Side (Standing): Viewing History Sent to Advertisers.} A tracking pixel that discloses video-viewing histories to a social-media platform resembles the tort of public disclosure of private facts.
\item \textbf{Opposite-Side (Standing): Injunction After a Breach.} Patients whose stolen records are threatened with release may seek forward-looking relief against an imminent risk under \textit{Clapper}.
\item \textbf{Fence-Sitter: Print Vendor.} Sending debt data to a mailing vendor whose staff never reads it: plaintiffs invoke disclosure torts; defendants say automated processing is not ``public'' (\textit{Barclift v. Keystone Credit Services} (3d Cir. 2024)).
\end{enumerate}

\section{Critique}

\textbf{Originalist critique.} The sharpest originalist attack came from Thomas: ``injury in fact'' appears nowhere in Article III, and the First Congress's copyright and penal statutes allowed suits without proof of damage. Scholars see tension with \textit{Vermont Agency v. Stevens} and \textit{SEC v. Jarkesy}, and call the analogue test a Burkean presumption against legislative novelty.

\textbf{Progressive critique.} Critics call the ``close relationship'' test indeterminate---why could defamation's falsity element be relaxed but not publication?---argue that suits between private parties pose no Article II problem, and warn that the decision pushes nationwide statutory class actions into state courts, undermining the Class Action Fairness Act.

\section{Key Quotations}

\begin{quote}
``Congress may enact legal prohibitions and obligations. And Congress may create causes of action for plaintiffs to sue defendants who violate those legal prohibitions or obligations. But under Article III, an injury in law is not an injury in fact. Only those plaintiffs who have been concretely harmed by a defendant's statutory violation may sue that private defendant over that violation in federal court.''
\end{quote}
\noindent --- Justice Kavanaugh, for the Court, 594 U.S. at 427.

\begin{quote}
``To have Article III standing to sue in federal court, plaintiffs must demonstrate, among other things, that they suffered a concrete harm. No concrete harm, no standing.''
\end{quote}
\noindent --- Justice Kavanaugh, for the Court, 594 U.S. at 417.

\begin{quote}
``In cases such as these where allegedly inaccurate or misleading information sits in a company database, the plaintiffs' harm is roughly the same, legally speaking, as if someone wrote a defamatory letter and then stored it in her desk drawer. A letter that is not sent does not harm anyone, no matter how insulting the letter is.''
\end{quote}
\noindent --- Justice Kavanaugh, for the Court, 594 U.S. at 434.

\end{document}
